\documentclass{article}
\usepackage{amsmath,amssymb,amsthm,algorithm,algorithmic}
\usepackage{fullpage}
\usepackage{enumitem}
\newtheorem{theorem}{Theorem}
\newtheorem{lemma}{Lemma}
\newtheorem{assumption}{Assumption}
\newcommand{\E}{\mathbb{E}}
\newcommand{\R}{\mathbb{R}}
\newcommand{\diag}{\operatorname{diag}}
\begin{document}

%%%%%%%%%%%%%%%%%%%%%%%%%%%%%%%%%%%%%%%%%%%%%%%%%%%%%
\section*{Randomized Coordinate Descent with Probability Weighting}
%%%%%%%%%%%%%%%%%%%%%%%%%%%%%%%%%%%%%%%%%%%%%%%%%%%%%

\subsection*{Mathematical Formulation}
Let $f:\R^{d}\to\R$ be a convex differentiable function.  
For each coordinate $i\in\{1,\dots ,d\}$ we are given a selection probability 
$p_i>0$ with $\sum_{i=1}^{d}p_i=1$.  

At iteration $t=0,1,2,\dots$ we draw a random index $i_t$ according to the
distribution $\Pr(i_t=i)=p_i$.  
Denote by $e_i\in\R^{d}$ the $i$‑th unit vector and by 
$\nabla_i f(w)=\partial f(w)/\partial w_i$ the $i$‑th partial derivative.  

The update rule is
\[
\boxed{
w_{t+1}=w_{t}-\eta\,\frac{1}{p_{i_t}}\,e_{i_t}\,\nabla_{i_t} f(w_t)
}
\tag{1}
\]
where $\eta>0$ is a (global) stepsize.  
The factor $1/p_{i_t}$ is crucial: it makes the stochastic direction an
unbiased estimator of the full gradient,
\[
\E\!\big[\,\tfrac{1}{p_{i_t}}e_{i_t}\nabla_{i_t} f(w_t)\mid w_t\big]
   =\sum_{i=1}^{d}p_i\frac{1}{p_i}e_i\nabla_i f(w_t)=\nabla f(w_t).
\tag{2}
\]

\subsection*{Pseudocode}
\begin{algorithm}[H]
\caption{Randomized Coordinate Descent (RCD)}\label{alg:RCD}
\begin{algorithmic}[1]
\REQUIRE Initial point $w_0\in\R^{d}$, stepsize $\eta>0$, probabilities $\{p_i\}_{i=1}^{d}$ with $\sum_i p_i=1$.
\FOR{$t=0,1,\dots,T-1$}
    \STATE Sample $i_t\sim\{1,\dots,d\}$ with $\Pr(i_t=i)=p_i$.
    \STATE Compute the partial gradient $g_{i_t}= \nabla_{i_t} f(w_t)$.
    \STATE Update the selected coordinate
    \[
        w_{t+1}= w_t - \eta\,\frac{1}{p_{i_t}}\,e_{i_t}\,g_{i_t}.
    \]
\ENDFOR
\RETURN $w_T$ (or the averaged iterate $\bar w_T=\frac1T\sum_{t=0}^{T-1}w_t$).
\end{algorithmic}
\end{algorithm}

\subsection*{Dry Run Example}
Consider the one–dimensional quadratic
\[
f(w)=(w-3)^2,\qquad w_0=0,\qquad \eta=0.1,
\]
and a single coordinate, hence $p_1=1$.  

\begin{center}
\begin{tabular}{c|c|c|c}
$t$ & $w_t$ & $\nabla f(w_t)=2(w_t-3)$ & $f(w_t)$\\\hline
0 & $0$ & $-6$ & $9$\\
1 & $w_1 = 0-\eta\frac{1}{1}(-6)=0+0.6=0.6$ &
$2(0.6-3)=-4.8$ & $(0.6-3)^2=5.76$\\
2 & $w_2 = 0.6-\eta(-4.8)=0.6+0.48=1.08$ &
$2(1.08-3)=-3.84$ & $(1.08-3)^2=3.6864$\\
3 & $w_3 = 1.08-\eta(-3.84)=1.08+0.384=1.464$ &
$2(1.464-3)=-3.072$ & $(1.464-3)^2=2.3666$\\
\end{tabular}
\end{center}
The iterate moves monotonically toward the minimizer $w^\star=3$, and the
objective value decreases at each step.

%%%%%%%%%%%%%%%%%%%%%%%%%%%%%%%%%%%%%%%%%%%%%%%%%%%%%
\section*{Assumptions}
%%%%%%%%%%%%%%%%%%%%%%%%%%%%%%%%%%%%%%%%%%%%%%%%%%%%%

\begin{assumption}[Convexity]\label{ass:convex}
$f$ is convex, i.e.
\[
f(y)\ge f(x)+\langle\nabla f(x),y-x\rangle,\qquad\forall x,y\in\R^{d}.
\]
\end{assumption}

\begin{assumption}[Coordinate‑wise $L_i$‑smoothness]\label{ass:smooth}
For each $i\in\{1,\dots,d\}$ there exists $L_i>0$ such that for all
$w\in\R^{d}$ and $\alpha\in\R$,
\[
\bigl|\nabla_i f(w+\alpha e_i)-\nabla_i f(w)\bigr|\le L_i|\alpha|.
\tag{3}
\]
Equivalently,
\[
f(w+\alpha e_i)\le f(w)+\alpha\nabla_i f(w)+\frac{L_i}{2}\alpha^{2}.
\tag{4}
\]
\end{assumption}

\begin{assumption}[Stepsize]\label{ass:stepsize}
The global stepsize satisfies
\[
0<\eta\le \frac{1}{\max_{i}L_i}.
\tag{5}
\]
\end{assumption}

\begin{assumption}[Existence of a minimizer]\label{ass:opt}
Let $w^\star\in\arg\min_{w\in\R^{d}}f(w)$ and define the weighted norm
\[
\|u\|_{P^{-1}}^{2}=u^{\top}P^{-1}u,\qquad 
P=\diag(p_{1},\dots,p_{d}).
\tag{6}
\]
We assume $\|w_{0}-w^\star\|_{P^{-1}}<\infty$.
\end{assumption}

%%%%%%%%%%%%%%%%%%%%%%%%%%%%%%%%%%%%%%%%%%%%%%%%%%%%%
\section*{Main Convergence Theorem}
%%%%%%%%%%%%%%%%%%%%%%%%%%%%%%%%%%%%%%%%%%%%%%%%%%%%%

\begin{theorem}[Expected $O(1/T)$ convergence]\label{thm:main}
Under Assumptions \ref{ass:convex}–\ref{ass:opt} and with the update
\eqref{1}, the iterates of Algorithm~\ref{alg:RCD} satisfy for any
$T\ge1$,
\[
\E\big[\,f(w_{T})-f(w^\star)\,\big]\le
\frac{\|w_{0}-w^\star\|_{P^{-1}}^{2}}{2\eta\,T}.
\tag{7}
\]
Consequently, the averaged iterate
$\bar w_T=\frac1T\sum_{t=0}^{T-1}w_t$ obeys the same bound because of convexity.
\end{theorem}

\subsection*{Theoretical Properties}
\begin{itemize}[leftmargin=*]
\item \textbf{Stability.} The condition \eqref{5} guarantees that each
coordinate update does not overshoot; the proof uses the smoothness
inequality \eqref{4}.
\item \textbf{Hyperparameter sensitivity.} The bound \eqref{7} is monotone
decreasing in $\eta$ until the maximal admissible value $1/\max_i L_i$.
Choosing $\eta$ close to this limit yields the smallest constant
$\frac{1}{2\eta}$.
\item \textbf{Comparison to full‑gradient GD.} For uniform probabilities
$p_i=1/d$, the weighted norm $\|\cdot\|_{P^{-1}}$ becomes $\sqrt{d}\,\|\cdot\|$,
and \eqref{7} recovers the standard $O(1/T)$ rate of gradient descent with
stepsize $1/L_{\max}$, up to the factor $\sqrt{d}$.
\end{itemize}

\subsection*{Discussion}
The theorem shows that, despite updating only a single coordinate per
iteration, the expected objective error decays as $1/T$, exactly as in
deterministic gradient descent for smooth convex problems. The key
ingredient is the unbiasedness \eqref{2} and the use of the weighted norm
$\|\cdot\|_{P^{-1}}$, which compensates for the non‑uniform sampling.

%%%%%%%%%%%%%%%%%%%%%%%%%%%%%%%%%%%%%%%%%%%%%%%%%%%%%
\section*{Preliminaries (Definitions and Lemmas)}
%%%%%%%%%%%%%%%%%%%%%%%%%%%%%%%%%%%%%%%%%%%%%%%%%%%%%

\begin{lemma}[One‑step descent]\label{lem:descent}
Let $w_{t+1}$ be generated by \eqref{1}. Under Assumption \ref{ass:smooth}
and stepsize condition \eqref{5},
\[
f(w_{t+1})\le f(w_t)-\frac{\eta}{2p_{i_t}}
\bigl(\nabla_{i_t} f(w_t)\bigr)^{2}.
\tag{8}
\]
\end{lemma}
\begin{proof}
From the smoothness inequality \eqref{4} with $\alpha=-\eta \nabla_{i_t}f(w_t)/p_{i_t}$,
\[
\begin{aligned}
f(w_{t+1})
&=f\!\bigl(w_t-\tfrac{\eta}{p_{i_t}}\nabla_{i_t}f(w_t) e_{i_t}\bigr)\\
&\le f(w_t)-\tfrac{\eta}{p_{i_t}}\bigl(\nabla_{i_t}f(w_t)\bigr)^2
   +\frac{L_{i_t}}{2}\Bigl(\tfrac{\eta}{p_{i_t}}\nabla_{i_t}f(w_t)\Bigr)^{2}.
\end{aligned}
\]
Because $\eta\le 1/L_{i_t}$ by \eqref{5},
\[
\frac{L_{i_t}}{2}\Bigl(\tfrac{\eta}{p_{i_t}}\nabla_{i_t}f(w_t)\Bigr)^{2}
\le\frac{\eta}{2p_{i_t}}\bigl(\nabla_{i_t}f(w_t)\bigr)^{2}.
\]
Subtracting yields \eqref{8}. ∎
\end{proof}

\begin{lemma}[Weighted distance recursion]\label{lem:dist}
Define $D_t:=\|w_t-w^\star\|_{P^{-1}}^{2}$. Then for any $t\ge0$,
\[
\boxed{
D_{t+1}\le D_t-2\eta\bigl(f(w_t)-f(w^\star)\bigr)+\eta^{2}
\sum_{i=1}^{d}\frac{(\nabla_i f(w_t))^{2}}{p_i}.
}
\tag{9}
\]
\end{lemma}
\begin{proof}
From the update \eqref{1},
\[
\begin{aligned}
w_{t+1}-w^\star
&=w_t-w^\star-\eta\frac{1}{p_{i_t}}e_{i_t}\nabla_{i_t}f(w_t).
\end{aligned}
\]
Taking the $P^{-1}$‑norm squared and using $P^{-1}e_i=e_i/p_i$,
\[
\begin{aligned}
D_{t+1}
&=(w_t-w^\star)^{\!\top}P^{-1}(w_t-w^\star)
   -\frac{2\eta}{p_{i_t}}(w_t-w^\star)^{\!\top}P^{-1}e_{i_t}\nabla_{i_t}f(w_t)\\
&\qquad+\frac{\eta^{2}}{p_{i_t}^{2}}e_{i_t}^{\!\top}P^{-1}e_{i_t}
      (\nabla_{i_t}f(w_t))^{2}.
\end{aligned}
\]
Since $P^{-1}e_{i}=e_i/p_i$, the middle term becomes
\[
-\frac{2\eta}{p_{i_t}}(w_t-w^\star)^{\!\top}\frac{e_{i_t}}{p_{i_t}}\nabla_{i_t}f(w_t)
=-\frac{2\eta}{p_{i_t}^{2}}(w_{t,i_t}-w^\star_{i_t})\nabla_{i_t}f(w_t).
\]
Similarly, the last term equals $\frac{\eta^{2}}{p_{i_t}^{3}}(\nabla_{i_t}f(w_t))^{2}$.
Taking expectation conditioned on $w_t$ and using $\Pr(i_t=i)=p_i$,
\[
\begin{aligned}
\E[D_{t+1}\mid w_t]
&=D_t-2\eta\sum_{i=1}^{d}\frac{1}{p_i}(w_{t,i}-w^\star_i)\nabla_i f(w_t)
   +\eta^{2}\sum_{i=1}^{d}\frac{(\nabla_i f(w_t))^{2}}{p_i}.
\end{aligned}
\tag{10}
\]
Convexity (Assumption \ref{ass:convex}) gives
\[
f(w_t)-f(w^\star)\le\langle\nabla f(w_t),w_t-w^\star\rangle
   =\sum_{i=1}^{d}(w_{t,i}-w^\star_i)\nabla_i f(w_t).
\]
Multiplying by $2\eta$ and substituting into \eqref{10} yields \eqref{9}. ∎
\end{proof}

%%%%%%%%%%%%%%%%%%%%%%%%%%%%%%%%%%%%%%%%%%%%%%%%%%%%%
\section*{Main Proof (Step‑by‑Step Derivation)}
%%%%%%%%%%%%%%%%%%%%%%%%%%%%%%%%%%%%%%%%%%%%%%%%%%%%%

We now prove Theorem~\ref{thm:main}. All expectations are taken over the
random choices $i_0,\dots,i_{T-1}$.

\begin{proof}
Let $D_t$ be as in Lemma~\ref{lem:dist}. Taking full expectation of
\eqref{9} and using the law of total expectation gives
\[
\E[D_{t+1}]\le \E[D_t]-2\eta\,\E\!\big[f(w_t)-f(w^\star)\big]
      +\eta^{2}\E\!\left[\sum_{i=1}^{d}\frac{(\nabla_i f(w_t))^{2}}{p_i}\right].
\tag{11}
\]
\textbf{[Step 1]}  Apply Lemma~\ref{lem:descent} and take expectation:
\[
\E\!\left[\sum_{i=1}^{d}\frac{(\nabla_i f(w_t))^{2}}{p_i}\right]
   =\E\!\left[\frac{1}{p_{i_t}}(\nabla_{i_t}f(w_t))^{2}\right]
   \le \frac{2}{\eta}\E\!\big[f(w_t)-f(w_{t+1})\big],
\tag{12}
\]
where the inequality follows from rearranging \eqref{8}.

\textbf{[Step 2]}  Substitute \eqref{12} into \eqref{11}:
\[
\E[D_{t+1}]
\le \E[D_t]-2\eta\,\E\!\big[f(w_t)-f(w^\star)\big]
   +2\eta\,\E\!\big[f(w_t)-f(w_{t+1})\big].
\tag{13}
\]

\textbf{[Step 3]}  Rearrange \eqref{13}:
\[
2\eta\,\E\!\big[f(w_t)-f(w^\star)\big]
\le \E[D_t]-\E[D_{t+1}]
   +2\eta\,\E\!\big[f(w_t)-f(w_{t+1})\big].
\tag{14}
\]

\textbf{[Step 4]}  Sum inequality \eqref{14} over $t=0,\dots,T-1$.
The telescoping sum $\sum_{t=0}^{T-1}(\E[D_t]-\E[D_{t+1}])$ equals
$\E[D_0]-\E[D_T]\ge 0$, and the terms
$\sum_{t=0}^{T-1}\E[f(w_t)-f(w_{t+1})]=\E[f(w_0)-f(w_T)]\le f(w_0)-f(w^\star)$.
Thus,
\[
2\eta\sum_{t=0}^{T-1}\E\!\big[f(w_t)-f(w^\star)\big]
\le \E[D_0] +2\eta\big(f(w_0)-f(w^\star)\big).
\tag{15}
\]

\textbf{[Step 5]}  Drop the non‑negative term $2\eta\big(f(w_0)-f(w^\star)\big)$ on the
right‑hand side of \eqref{15} to obtain a simpler bound:
\[
2\eta\sum_{t=0}^{T-1}\E\!\big[f(w_t)-f(w^\star)\big]\le \E[D_0].
\tag{16}
\]

\textbf{[Step 6]}  Divide both sides of \eqref{16} by $2\eta T$ and use convexity
of $f$ to relate the average of the function values to the function value at
the averaged iterate (or directly to $f(w_T)$, both satisfy the same bound):
\[
\frac{1}{T}\sum_{t=0}^{T-1}\E\!\big[f(w_t)-f(w^\star)\big]
\le \frac{\E[D_0]}{2\eta T}
   =\frac{\|w_0-w^\star\|_{P^{-1}}^{2}}{2\eta T}.
\tag{17}
\]

Since $f(w_T)-f(w^\star)$ is one element of the average, it follows that
\[
\E\big[f(w_T)-f(w^\star)\big]\le\frac{\|w_0-w^\star\|_{P^{-1}}^{2}}{2\eta T},
\]
which is exactly the claim \eqref{7}.  ∎
\end{proof}

%%%%%%%%%%%%%%%%%%%%%%%%%%%%%%%%%%%%%%%%%%%%%%%%%%%%%
\section*{Rate Derivation (With Constants)}
%%%%%%%%%%%%%%%%%%%%%%%%%%%%%%%%%%%%%%%%%%%%%%%%%%%%%

The bound \eqref{7} can be rewritten using
$L_{\max}:=\max_i L_i$ and the uniform probability case $p_i=1/d$:
\[
\E\big[f(w_T)-f(w^\star)\big]\le
\frac{d\,\|w_0-w^\star\|^{2}}{2\eta T}
\le\frac{d\,L_{\max}\|w_0-w^\star\|^{2}}{2T},
\]
where we used the admissible stepsize $\eta=1/L_{\max}$.  Hence the
asymptotic rate is $O(d/T)$ in the worst case, but the dependence on $d$
is mitigated when the sampling distribution matches the coordinate
smoothness, e.g. $p_i\propto L_i$.

%%%%%%%%%%%%%%%%%%%%%%%%%%%%%%%%%%%%%%%%%%%%%%%%%%%%%
\section*{Special Cases}
%%%%%%%%%%%%%%%%%%%%%%%%%%%%%%%%%%%%%%%%%%%%%%%%%%%%%

\begin{itemize}[leftmargin=*]
\item \textbf{Uniform sampling.} $p_i=1/d$ for all $i$.  The weighted norm
$\|\,\cdot\,\|_{P^{-1}}$ reduces to $\sqrt{d}\,\|\cdot\|$, giving the classic
$O(d/T)$ bound.

\item \textbf{Importance sampling.} Choose $p_i = L_i/\sum_j L_j$.
Then $\sum_i (\nabla_i f)^2/p_i = (\sum_j L_j)\sum_i (\nabla_i f)^2/L_i$,
which can be substantially smaller than the uniform case, yielding a
tighter constant in \eqref{7}.

\item \textbf{Strongly convex $f$.} If $f$ is $\mu$‑strongly convex,
the same analysis combined with the PL inequality yields a linear rate
$\E[f(w_T)-f(w^\star)]\le (1-\eta\mu)^T\big(f(w_0)-f(w^\star)\big)$.
\end{itemize}

%%%%%%%%%%%%%%%%%%%%%%%%%%%%%%%%%%%%%%%%%%%%%%%%%%%%%
\end{document}